\begin{afterword}
\summaryhead

The essay asks how much evidence a belief requires and answers with a lottery. With seventy balls and six to match there are 131,115,985 possible tickets. A box that always beeps for the winning ticket, and beeps for a losing one a quarter of the time, has a likelihood ratio of four to one. It cuts the field by a factor of four but leaves over thirty million losing tickets that also beep. Independent boxes multiply.

Evidence as strong as the number of alternatives gives only even odds, because about one loser will pass along with the winner. Measured in bits, each box adds two; twenty-eight bits give odds of about two to one, and thirty-four bits give better than 99 per cent confidence. In general, a larger space of possibilities, a less likely hypothesis or a wish for more confidence all require more evidence. Believing without it, the essay says, is like driving without fuel.

\responsehead

The title asks a general question. The essay answers a special one: how much evidence it takes to pick one item from a list of known length, all equally likely, using tests that never miss the true answer and never interfere with each other. In that case the answer is arithmetic. The essay names the factor that decides every real case, how unlikely the hypothesis ``seems a priori,'' in one clause of its general rule, and says nothing about how to set it. Outside a lottery, the honest answer to the title is that it depends on a number you usually do not have. The essay does not give that answer.

The arithmetic it does give is careless. Odds are rounded without notice. Probability and odds, which the essay takes a sentence to distinguish, are confused three times. The bit is defined as the information in an outcome and then used as the logarithm of a likelihood ratio, a different quantity that happens to coincide in the lottery. In the paragraph that turns the calculation into a moral, ten boxes are said to let through 131 losing tickets for each winner; the number is 125. Two sentences later: ``That's not a pointless bureaucratic regulation; it's math.''

The moral itself is misstated. ``You cannot form accurate beliefs based on inadequate evidence'' is false; a guess can be right. What thin evidence cannot give is warranted confidence. The fuel analogy needs the false version, since a car without fuel cannot move at all.

The unit it leaves the reader with, bits, sounds like a measurement, but the essay never shows how to compute it for any belief outside a lottery. The idea of measuring evidence on a log scale belongs to Turing and I. J. Good, and neither is named.

In short: it answers its title only for a lottery, gets the lottery's arithmetic wrong at the moment it says ``it's math,'' and leaves the reader with a unit the essay never shows how to apply to anything they believe.
\end{afterword}
